\chapter{问题背景、任务重述与分析}

\section{工程背景}
FAST 主动反射面系统通过促动器驱动主索节点，实现从基准球面到工作抛物面的动态调节。与传统固定口径望远镜相比，主动反射面具有更强的目标适应能力，但也引入了复杂的结构约束与几何耦合关系。建模的核心目标是在可执行约束下最大化电磁波聚焦效果。

\section{问题重述}
结合题面，可将任务概括为三步：
\begin{enumerate}
  \item 在给定观测方向时确定理想抛物面几何参数；
  \item 建立节点调节与促动器行程之间的约束优化模型，使工作面尽量贴合理想面；
  \item 根据调节结果计算馈源舱接收比，并与基准球面接收比进行比较。
\end{enumerate}

\section{问题难点}
\begin{enumerate}
  \item \textbf{规模大}：节点数与连接边数较多，变量维度高；
  \item \textbf{约束强}：促动器行程和相邻边长变化均有严格边界；
  \item \textbf{目标间接}：反射接收性能并非直接优化变量，需要二次光学评估。
\end{enumerate}

\section{建模思路}
本文采用“几何拟合 + 结构约束 + 光学验证”的三段式建模路线：
\begin{enumerate}
  \item 用理想抛物面确定理论目标形状；
  \item 用约束优化求解节点位移与促动器量；
  \item 用光线追踪评估接收比，验证调节效果。
\end{enumerate}

\section{本章小结}
本章明确了题目要求、工程意义与技术难点。后续章节将分别给出数学模型、求解算法和结果分析。
